\section{Tensor identities and the scope of further transfers}\label{sec:functorial-identities}

This appendix records exact identities which distinguish the exterior-square argument used in the main theorem from stronger functorial transfers that would settle the rank-five comparison directly. Write $V=\rho_4$, $R=\rh$, and $W=\rho_6$.

\begin{proposition}[Finite identities]\label{prop:finite-tensors}
The following identities hold as complex $S_5$-representations:
\begin{align*}
R\otimes R&=\one\oplus V\oplus V\eta\oplus R\oplus R\eta\oplus W,\\
V\otimes R&=V\oplus R\oplus R\eta\oplus W,\\
W\otimes R&=V\oplus V\eta\oplus R\oplus R\eta\oplus W^{\oplus2},\\
\operatorname{Sym}^2V&=\one\oplus V\oplus R,\\
\operatorname{Sym}^2R&=\one\oplus V\oplus R\oplus R\eta,\qquad
\wedge^2R=V\eta\oplus W.
\end{align*}
\end{proposition}

\begin{proof}
For a permutation $g$, let $a$ and $b$ be its numbers of fixed letters and two-cycles. The characters of $V$ and $R$ are $a-1$ and $a(a-3)/2+b$, respectively: the permutation character on unordered pairs is $\binom a2+b$, and subtracting the five-letter permutation character gives $R$. The character of $W$ is
\[
\chi_W(g)=\frac{\chi_V(g)^2-\chi_V(g^2)}2.
\]
Use the class order in \eqref{eq:rhochar}, with sizes $1,10,15,20,30,24,20$. Multiplying characters gives the first three rows; using
\[
\chi_{\operatorname{Sym}^2 X}(g)=\frac{\chi_X(g)^2+\chi_X(g^2)}2,
\qquad
\chi_{\wedge^2 X}(g)=\frac{\chi_X(g)^2-\chi_X(g^2)}2
\]
gives the last two. Direct substitution verifies equality on every class, which proves equality of the semisimple representations. For example the symmetric square of $V$ has character $(10,4,2,1,0,0,1)$, the character of the permutation action on unordered pairs. Thus it is exactly $\one\oplus V\oplus R$.
\end{proof}

\begin{proposition}[Exceptional tensor comparisons]\label{prop:tensor-cancellation}
At an exceptional prime put $T=p^{-s}$ and $\Phi_p=((1+T)/(1-T))^2$. The local quotients of the self tensor, symmetric square, and exterior square of the automorphic rank-five parameter by their Artin counterparts are respectively
\[
\Phi_p^2,\qquad\Phi_p,\qquad\Phi_p.
\]
Tensoring that parameter with either $V$ or $W$ instead gives quotient one. Twisting the original rank-five comparison by $\eta$ gives $\Phi_p^{-1}$.
\end{proposition}

\begin{proof}
At a four-cycle the spectra of $V$ and $W$ are
\[
\{1,i,-1,-i\},\qquad\{1,-1,i,i,-i,-i\}.
\]
The two rank-five spectra are $r=\{1,-1,-1,i,-i\}$ and $a=\{1,1,1,i,-i\}$. Counting products, with coordinates ordered by the eigenvalues $1,i,-1,-i$, gives
\[
\begin{array}{c|rrrr|rrrr}
 &\multicolumn{4}{c|}{\text{Artin}}&\multicolumn{4}{c}{\text{automorphic}}\\
\text{self tensor}&7&6&6&6&11&6&2&6\\
\operatorname{Sym}^2&5&3&4&3&7&3&2&3\\
\wedge^2&2&3&2&3&4&3&0&3\\
V\text{ tensor}&5&5&5&5&5&5&5&5\\
W\text{ tensor}&7&8&7&8&7&8&7&8.
\end{array}
\]
Taking reciprocal Euler polynomials proves each assertion. Finally $\eta(p)=-1$, so twisting replaces $T$ by $-T$ and inverts $\Phi_p$.
\end{proof}

In particular the two rank-five terms in the Artin decomposition of $R\otimes R$ have comparison quotients $\Phi_p$ and $\Phi_p^{-1}$, which cancel. The remaining self Rankin--Selberg quotient is $\Phi_p^2$; the holomorphy of its numerator does not exclude poles of the quotient at zeros of its denominator. By contrast, $\wedge^2R$ contains neither rank-five term, and its intrinsically entire Hecke function is the separate function $C$ used in Proposition~\ref{prop:commonpole}.

\begin{proposition}[Two conditional transfer equivalences]\label{prop:transfer-equivalence}
Let $\pi_4$ be the matched cuspidal representation for $V$, and let $\tau$ be the matched cuspidal representation over $E$ for either faithful three-dimensional $A_5$-representation $U=\operatorname{Ad}\varrho$. Each of the following is equivalent to the existence of a unitary cuspidal weak automorphic realization of $R$:
\begin{enumerate}[label=(\roman*)]
\item a unitary isobaric weak symmetric-square transfer of $\pi_4$ to $\GL_{10}/\Q$;
\item a unitary isobaric weak $\operatorname{Asai}^+$ transfer of $\tau$ to $\GL_9/\Q$.
\end{enumerate}
Here weak means equality of unramified local parameters outside a finite set. The plus convention for tensor induction is the one for which an odd element $g$ has trace $\chi_U(g^2)$.
\end{proposition}

\begin{proof}
The Rankin--Selberg inputs below are the unequal-rank boundary regularity of \citet[Proposition~3.3, p.~795]{jacquetshalika1981ii}, their equal-rank simple-pole criterion \citep[Proposition~3.6, pp.~802--804]{jacquetshalika1981ii}, and nonvanishing on the line one \citep[Theorem~3.7, p.~804]{jacquetshalika1981ii}. The last input is Shahidi's theorem, also stated directly in \citet[p.~462]{shahidi1980}. These results concern partial Rankin--Selberg functions of unitary cuspidal representations; applying them constituent by constituent gives the stated pole multiplicities for unitary isobaric sums.

The finite identity for (i) is $\operatorname{Sym}^2V=\one\oplus V\oplus R$. For (ii), the two three-dimensional characters of $A_5$ have values $3,-1,0,(1+\sqrt5)/2,(1-\sqrt5)/2$, with the last two interchanged. Their product is the character of $V|_{A_5}\oplus R|_{A_5}$. At odd elements of types $2,4,3\cdot2$, the plus tensor-induction traces are respectively $3,-1,0$, because their squares have types $1,2^2,3$. These are exactly the traces of $V\oplus R$. Therefore
\[
\operatorname{Asai}_{A_5}^{S_5,+}U=V\oplus R,
\qquad
\operatorname{Asai}_{A_5}^{S_5,-}U=V\eta\oplus R\eta.
\]

Suppose first that the transfer in (i) exists, and call it $\Sigma$. All of its good Satake eigenvalues have absolute value one, since they are the stated finite-image eigenvalues. In the Euler logarithm of $L^S(s,\Sigma\times\pi_4^\vee)$, higher prime powers are therefore bounded as real $s\downarrow1$. Chebotarev gives prime-term coefficient
\[
\langle\chi_{\one+V+R},\chi_V\rangle_{S_5}=1
\]
in front of $\log(1/(s-1))$. Rankin--Selberg pole theory and nonvanishing on the line one show that $\pi_4$ occurs with multiplicity one in the unitary isobaric decomposition of $\Sigma$. The same argument with the trivial character shows that $\one$ occurs once. Removing those two constituents leaves a unitary isobaric representation $\Sigma_5$ of rank five whose good parameters are exactly those of $R$.

If $\Sigma_5=\boxplus_j\sigma_j^{\boxplus m_j}$ with distinct unitary cuspidal $\sigma_j$, the pole order of its self Rankin--Selberg function is $\sum_jm_j^2$. The finite character computation instead gives
\[
\langle\chi_R,\chi_R\rangle_{S_5}=1.
\]
Thus there is exactly one constituent of multiplicity one, proving cuspidality. This proves the forward implication in (i). In (ii), the same argument extracts $\pi_4$ once from the rank-nine transfer and proves that its rank-five complement is cuspidal. Conversely a weak cuspidal realization $\sigma_5$ of $R$ gives the respective transfers
\[
\one\boxplus\pi_4\boxplus\sigma_5,
\qquad
\pi_4\boxplus\sigma_5.
\]
These have exactly the required good parameters by the finite identities.

A weak cuspidal realization $\sigma_5$ also settles the specified comparison with $\pi$. Their good parameters differ only at a finite set and at $\Omega$, and $\sum_{p\in\Omega}p^{-1}<\infty$. Their Euler logarithms therefore show that
\[
\log L^S(s,\pi\times\sigma_5^\vee)
-\log L^S(s,\sigma_5\times\sigma_5^\vee)
\]
remains bounded as real $s\downarrow1$; all higher prime powers are uniformly summable. The second function has a simple pole, so the first does as well. Rankin--Selberg pole theory gives $\pi\simeq\sigma_5$. The good discrepancy set is then finite, and Proposition~\ref{prop:package}(d) gives equality of the standard local factors at every place.
\end{proof}

These equivalences specify an additional transfer that would solve the target; they do not supply it. \citet[Proposition 3.1(b)]{asgariraghuram} explicitly distinguish the conjectural symmetric-square transfer from the established exterior-square transfer in the tensor-product setting. Over $E$, the particular tensor product of the two adjoint three-dimensional representations is already an isobaric sum of the known four- and five-dimensional $A_5$ constituents. Descending it with the correct inert-prime parameters is the same unresolved choice encountered in the rank-five construction. The analytic exterior-square holomorphy used in the main proof is available without making either additional transfer assumption.

\subsection{A symmetric-square bound subsumed by the Artin-factor criterion}

Symmetric-square holomorphy gives a further Dedekind-zeta zero at a comparison pole. The factorization below shows exactly why this does not strengthen the earlier Artin-factor variant in Remark~\ref{rem:weaker}.

\begin{proposition}[Symmetric-square condition and its redundancy]\label{prop:symmetric-square-redundancy}
Let $D_8=C_{S_5}((12)(34))$, let $T=L^{D_8}$, and put $D(s)=\zeta_T(s)$. Then
\[
\operatorname{Sym}^2R=\Ind_{D_8}^{S_5}\one,\qquad [T:\Q]=15,
\]
and $T$ has signature $(3,6)$ and infinite factor $\Gr(s)^9\Gr(s+1)^6$.
At every pole $z$ of $M$ of order $r>0$ in $0<\Rea z<1$,
\[
r\leq\min\{\ord_zB,\ord_zC,\ord_zD\}.
\]
Thus absence of a common off-central real zero of $B,C,D$ suffices for the all-place and entire-automorphy conclusions of Theorem~\ref{thm:main}, with no restriction on central multiplicities. However,
\begin{equation}\label{eq:symmetric-square-redundancy}
D(s)=L(s,V)L(s,R)\zeta_F(s),\qquad
\ord_zD\geq\ord_zL(s,R)\quad(0<\Rea z<1).
\end{equation}
The $D$ pole bound is consequently already implied by the prior $L(s,R)$ bound; the three-function condition is subsumed by Remark~\ref{rem:weaker}.
\end{proposition}

\begin{proof}
The centralizer of $(12)(34)$ fixes 5, independently interchanges the two letters in each of its two orbits, and may interchange the orbits. It has order $2^2\cdot2=8$ and is dihedral, generated by $(1324)$ and $(12)$. Its cosets identify equivariantly with the fifteen double transpositions, by $xD_8\mapsto x(12)(34)x^{-1}$. A fixed point of $g$ in this action is exactly a double transposition commuting with $g$.

The symmetric-square character formula in Proposition~\ref{prop:finite-tensors} gives the full calculation
\[
\begin{array}{c|rrrrrrr}
\text{cycle type}&1&2&2^2&3&4&5&3\cdot2\\ \hline
\chi_R(g^2)&5&5&5&-1&1&0&-1\\
\chi_{\operatorname{Sym}^2R}(g)&15&3&3&0&1&0&0\\
\text{fixed double transpositions}&15&3&3&0&1&0&0.
\end{array}
\]
For completeness, $(12)$ commutes with exactly $(12)(34),(12)(35),(12)(45)$. The centralizer of $(12)(34)$ contains exactly the three double transpositions on $\{1,2,3,4\}$. A three-cycle or a permutation of type $3\cdot2$ has centralizer $C_3\times C_2$, containing no double transposition. A four-cycle has cyclic centralizer of order four, whose square is its unique double transposition. A five-cycle has cyclic centralizer of order five and contributes none. These counts prove the last row and hence the representation identity on every class. Artin induction gives $L(s,\operatorname{Sym}^2R)=\zeta_T(s)$, including ramified factors. Complex conjugation fixes three cosets and exchanges the remaining twelve in six pairs, giving signature $(3,6)$. Finally $\Gc(s)=\Gr(s)\Gr(s+1)$ gives the stated infinite factor.

At an exceptional prime the two rank-five spectra in Proposition~\ref{prop:package} have symmetric-square multiplicities
\[
\begin{array}{c|rrrr}
\text{eigenvalue}&1&i&-1&-i\\ \hline
\text{Artin}&5&3&4&3\\
\text{automorphic}&7&3&2&3 .
\end{array}
\]
These count products with indices $i\leq j$, including diagonal products. With $x=p^{-s}$ the two local factors are
\[
D_p=(1-x)^{-5}(1+x)^{-4}(1+x^2)^{-3},\qquad
\mathcal H_p=(1-x)^{-7}(1+x)^{-2}(1+x^2)^{-3}.
\]
Therefore $\mathcal H_p/D_p=((1+x)/(1-x))^2$, with the same quotient direction as $M$. At every other undeleted prime the parameters match.

The external analytic input is \citet[Theorem~7.1, pp.~49--50]{takeda2015}: for a cuspidal representation with unitary central character and a unitary character $\chi$, the partial twisted symmetric-square function is holomorphic except for possible simple poles at 0 and 1, permitted only if $\chi^r\omega_\pi^2=1$. The deleted set can consist exactly of the infinite places, those over 2, and the places where $\pi$ or $\chi$ ramifies. Apply this with rank five and $\chi=1$. In particular it supplies holomorphy throughout the open strip. We use this analytic theorem as a published input; it neither asserts nonvanishing there nor constructs a $\GL_{15}$ transfer.

Take a finite set $S$ containing those places, all primes ramified in $L$, and the original finite exceptions. Enlarging Takeda's set deletes unramified factors by multiplication by entire polynomials $\prod_{i\leq j}(1-\alpha_{p,i}\alpha_{p,j}p^{-s})$, so
\[
\mathcal H^S(s)=L^S(s,\pi,\operatorname{Sym}^2)
\]
remains holomorphic in the strip. These added polynomials are nonzero for $\Rea s>0$, since $\pi$ is tempered. On the Artin side every local eigenvalue on inertia invariants is a root of unity, or the invariant space is absent. Thus $D^S/D$ is a finite product of polynomials $1-\lambda p^{-s}$ with $|\lambda|=1$, and is holomorphic and nowhere zero for $\Rea s>0$. The finitely removed factors of $Q_\Omega$ likewise have no zero or pole there, since $1\pm p^{-s}\ne0$.

Absolute convergence in $\Rea s>1$ and the local calculation give
\[
\mathcal H^S(s)=D^S(s)Q_{\Omega\setminus S}(s).
\]
The equality continues meromorphically. The established completed comparison and the preceding deletion factors therefore give, on the open strip,
\[
M(s)=u(s)\frac{\mathcal H^S(s)}{D(s)}
\]
with $u$ holomorphic and nowhere zero. At a pole of order $r$ this implies
\[
\ord_zD=r+\ord_z\mathcal H^S\geq r.
\]
Together with Proposition~\ref{prop:commonpole} this proves the three order bounds at every complex point in the strip. The pole of $\zeta_T$ at 1 is outside the domain. Excluding off-central real triple zeros excludes all such real poles of $M$, and the same central-divisor argument used in Theorem~\ref{thm:main} gives its conclusions.

To establish the precise limitation, Proposition~\ref{prop:finite-tensors} and \eqref{eq:factorization} give
\[
\zeta_T=\zeta_{\Q}L(V)L(R)L(R\eta)=L(V)L(R)\zeta_F.
\]
The functions $L(V)$ and $\zeta_F$ are holomorphic in the strip, proving \eqref{eq:symmetric-square-redundancy}. The existing standard comparison already gives $\ord_zL(s,R)=r+\ord_zL(s,\pi)\geq r$ at a pole of $M$. Thus it proves the $D$ bound without the symmetric-square analytic theorem.

More explicitly, write $Z(f)$ for the set of points of positive meromorphic order in the open strip. Since $B=L(\eta)L(R)$ with $L(\eta)$ entire, the exact inclusions are
\[
Z(L(R))\cap Z(C)
\ \subseteq\ Z(B)\cap Z(C)\cap Z(D)
\ \subseteq\ Z(B)\cap Z(C).
\]
They remain valid after restriction to off-central real points. Hence the triple-zero exclusion is weaker than the $B,C$-only formulation but at least as restrictive as the already available $L(R),C$ formulation. No new criterion beyond that Artin-factor variant follows.
\end{proof}

\begin{remark}[Formal and arithmetic scope]
The displayed identities and holomorphy constraints alone do not imply that every common $B,C$ zero is a $D$ zero. As a local-germ example, assign $L(\eta)=L(V\eta)=t$ and $L(R)=L(R\eta)=L(V)=L(W)=\zeta_{\Q}=\zeta_F=M=L(\pi)=L(\pi\eta)=\mathcal H^S=1$. Then $B=C=t$ and $D=1$, consistently with these identities. This is only a formal example, not an $S_5$ field or a global system of Euler products. Any arithmetic common $B,C$ zero missed by $D$ must be a zero of $L(\eta)$, by \eqref{eq:symmetric-square-redundancy}. No strict arithmetic zero-set inclusion or newly certified field is established.
\end{remark}
